\chapter{Literature Review}
\label{ch:lit}

% Helper: author/venue line under each paper heading
\newcommand{\paperinfo}[1]{\noindent\textit{#1}\par\smallskip}

\section{Introduction}

Recent research covers three areas that are relevant to this work: autonomous
attackers (red), autonomous defenders (blue), and reasoning and learning methods.
The works reviewed below show that each area is advancing quickly. However, no
single system joins them into one self-improving loop. All 25 papers were read
and summarised from their abstracts. They are grouped into six themes, starting
with the works closest to PurpleAgents.

% -------------------------------------------------------------------
\section{Autonomous and Agentic Purple Teaming}

\subsection{Toward an Autonomous Purple Teaming Framework for Security and Safety in LLMs (AutoPT)}
\paperinfo{L. H. Park et al., PRDC 2025~\cite{park2025autopt}}
This work runs two co-evolving agents in a loop. A red module uses
coverage-guided fuzzing to find jailbreaks, and a blue module self-heals its
guardrails using that feedback. It is the closest published match to our idea.
\textbf{Gap:} it covers LLM safety only and relies on a single measurement score,
with no dual reasoning.

\subsection{Purple-teaming LLMs with Adversarial Defender Training (PAD)}
\paperinfo{A. Zhou et al., arXiv:2407.01850, 2024~\cite{zhou2024pad}}
PAD trains an attacker and a defender together, so that the defender learns to
catch and explain unsafe responses. We cite it because it shows that red/blue
co-training works. \textbf{Gap:} it focuses on guardrail training only and does
not form a closed loop.

\subsection{From Promise to Peril: Rethinking Red and Blue Teaming in the Age of LLMs}
\paperinfo{A. Abuadbba et al., arXiv:2506.13434, 2025~\cite{abuadbba2025peril}}
This paper maps the use of LLMs in red and blue teaming onto known security
frameworks and points out weaknesses such as hallucination. It supports our
choice to keep humans in the loop. \textbf{Gap:} it is a position paper; no
system is built.

% -------------------------------------------------------------------
\section{Offensive (Red) Agents}

\subsection{VulnBot: Autonomous Penetration Testing for a Multi-Agent Collaborative Framework}
\paperinfo{H. Kong et al., arXiv:2501.13411, 2025~\cite{kong2025vulnbot}}
VulnBot is a multi-agent LLM system that splits a penetration test into
reconnaissance, scanning, and exploitation. It gives us a template for
structuring our red agents. \textbf{Gap:} it covers attack only, with no defense
or learning.

\subsection{RedTeamLLM: An Agentic AI Framework for Offensive Security}
\paperinfo{B. Challita and P. Parrend, arXiv:2505.06913 / IFIP AI4S 2025~\cite{challita2025redteamllm}}
RedTeamLLM is an agentic penetration-testing system built on a
summarize--reason--act loop. It addresses plan correction and memory. It informs
the reasoning and memory design of our red side. \textbf{Gap:} it is offensive
only, with no validation step and no memory of lessons learned.

\subsection{Autonomous Adversary: Red-Teaming in the Age of LLM}
\paperinfo{A. K. M. Mamun et al., arXiv:2605.06486, 2026~\cite{mamun2026adversary}}
This work benchmarks LLM-driven attacks and uses an LLM as a judge to check
whether each attack step succeeded. It informs how our validation step decides
success. \textbf{Gap:} it is a red-side benchmark with no defense loop.

\subsection{PentestGPT: Evaluating and Harnessing LLMs for Automated Penetration Testing}
\paperinfo{G. Deng et al., USENIX Security 2024, arXiv:2308.06782~\cite{deng2024pentestgpt}}
PentestGPT splits the work into three cooperating modules so that the model does
not lose track of the test. It is a foundation for our modular red agent.
\textbf{Gap:} it is a penetration-testing tool with single-path reasoning and no
blue side.

\subsection{AutoAttacker: An LLM-Guided System to Implement Automatic Cyber-attacks}
\paperinfo{J. Xu et al., arXiv:2403.01038, 2024~\cite{xu2024autoattacker}}
AutoAttacker focuses on the post-breach, ``hands-on-keyboard'' stage of an
attack. It shows the kind of post-exploitation actions a red agent may need to
simulate. \textbf{Gap:} it covers post-breach attacks only, with no defense.

\subsection{LLM Agents Can Autonomously Hack Websites}
\paperinfo{R. Fang et al., arXiv:2402.06664, 2024~\cite{fang2024hack}}
This work shows that agents built on a leading LLM can perform SQL injection
without prior knowledge of the flaw, with a reported success rate of about 73\%.
It is concrete proof that autonomous attackers are real. \textbf{Gap:} it is an
attack demonstration with no defense side.

\subsection{On the Feasibility of Using LLMs to Execute Multistage Network Attacks (Incalmo)}
\paperinfo{B. Singer et al., arXiv:2501.16466, 2025~\cite{singer2025incalmo}}
This paper finds that plain LLMs fail at multi-host attacks until an abstraction
layer is added (an attack-graph service, an environment-state service, and an
action planner). These services resemble our attack-graph reasoning and evidence
layer. \textbf{Gap:} it addresses attack execution only, with no reasoning or
learning loop.

% -------------------------------------------------------------------
\section{Defensive and SOC Agents}

\subsection{CORTEX: Collaborative LLM Agents for High-Stakes Alert Triage}
\paperinfo{B. Wei et al., arXiv:2510.00311, 2025~\cite{wei2025cortex}}
In CORTEX, specialised agents collaborate over real evidence: a
behaviour-analysis agent, evidence-gathering agents, and a reasoning agent that
makes an auditable decision. It is the closest match to our blue side and our
evidence-based reasoning. \textbf{Gap:} it performs blue-side triage only, with no
red side and no loop.

\subsection{Large Language Models for Security Operations Centers: A Comprehensive Survey}
\paperinfo{arXiv:2509.10858, 2025~\cite{soc2025survey}}
This survey explains how LLMs fit into the security operations centre (SOC)
workflow, including triage, correlation, and prioritisation. It grounds what our
defensive agents are expected to do. \textbf{Gap:} it is a survey; no system is
built.

\subsection{A Survey of Agentic AI and Cybersecurity: Challenges, Opportunities and Use-case Prototypes}
\paperinfo{S. J. Lazer et al., arXiv:2601.05293, 2026~\cite{lazer2026agentic}}
This survey studies agentic AI in security and warns of agent-specific risks
such as memory poisoning. It also includes prototypes. It supports our
multi-agent design and our human-approval layer. \textbf{Gap:} it offers a survey
and prototypes, but no closed loop.

% -------------------------------------------------------------------
\section{Reinforcement-Learning Cyber Operations}

\subsection{CybORG: An Autonomous Cyber Operations Research Gym}
\paperinfo{C. Baillie et al., arXiv:2002.10667, 2020~\cite{baillie2020cyborg}}
CybORG is an open simulation ``gym'' in which red (attacker) and blue (defender)
agents are trained together using reinforcement learning. It shows the kind of
safe range needed for co-training agents. \textbf{Gap:} it is a training
environment only, with no reasoning engine.

\subsection{Deep Reinforcement Learning for Autonomous Cyber Operations: A Survey}
\paperinfo{C. Palmer et al., arXiv:2310.07745, 2024~\cite{palmer2024drl}}
This survey frames cyber defense as a game between a blue defender and a red
attacker and lists the main open challenges. It supports our training setup.
\textbf{Gap:} it is an RL survey with no dual reasoning.

\subsection{Entity-based Reinforcement Learning for Autonomous Cyber Defence}
\paperinfo{I. Symes Thompson et al., arXiv:2410.17647, 2024~\cite{symes2024entity}}
This work builds a defensive agent that keeps working when the network changes,
by viewing the network as a set of entities. It guides how our defensive agents
can stay robust across different systems. \textbf{Gap:} it covers defense only and
uses single-path decisions.

\subsection{Hierarchical Multi-agent Reinforcement Learning for Cyber Network Defense}
\paperinfo{A. Singh et al., arXiv:2410.17351, AAMAS 2025~\cite{singh2025hmarl}}
This paper splits defense into sub-tasks handled by several agents under a master
policy. It is a model for how our blue agents can cooperate. \textbf{Gap:} it
addresses blue-side coordination only, with no reasoning layer.

\subsection{Towards Autonomous Cyber Operation Agents: Exploring the Red Case}
\paperinfo{L. Li et al., arXiv:2309.02247, 2023~\cite{li2023redcase}}
This work trains an attacking red agent with reinforcement learning and studies
how to make it work across different networks. It supports the building of our
offensive agents. \textbf{Gap:} it covers red-agent training only.

\subsection{A Red Teaming Framework for Evaluating Robustness of AI-enabled SOAR Systems}
\paperinfo{A. J. Shaikh et al., arXiv:2605.17075, 2026~\cite{shaikh2026soar}}
This framework combines an LLM planner with a reinforcement-learning controller
to run multi-stage attacks against automated defenders. It finds that an LLM
alone cannot sustain multi-stage attacks. It is a real example of an
attacker-versus-defender loop. \textbf{Gap:} it focuses on attack testing only,
with no learning loop.

% -------------------------------------------------------------------
\section{Dual and Abductive Reasoning}

\subsection{Evidential Cyber Threat Hunting}
\paperinfo{F. Araujo et al., SIAM AI4CS 2021, arXiv:2104.10319~\cite{araujo2021evidential}}
This paper proposes a formal reasoning framework that works over three spaces
--- knowledge, hypothesis, and action --- and attaches a confidence level to each
conclusion. It is the academic basis of our dual-reasoning engine. \textbf{Gap:} it
is a reasoning framework only and is not integrated into an attack--defense
system.

\subsection{Diagnosing Unknown Attacks in Smart Homes Using Abductive Reasoning}
\paperinfo{K. Ramkumar et al., IEEE, arXiv:2412.10738, 2025~\cite{ramkumar2025abductive}}
This work combines anomaly detection with abductive reasoning to find the class
of an unknown attack and the requirement it violated. It gives a concrete method
for our hypothesis-first path. \textbf{Gap:} it performs diagnosis only, with no
attack--defense loop.

% -------------------------------------------------------------------
\section{Surveys and Foundations}

\subsection{When LLMs Meet Cybersecurity: A Systematic Literature Review}
\paperinfo{J. Zhang et al., \textit{Cybersecurity} (Springer), 2025, arXiv:2405.03644~\cite{zhang2025llmcyber}}
This review covers more than 300 works across 25 LLMs and over 10 security tasks.
It is the main peer-reviewed foundation for our related work. \textbf{Gap:} it is
a broad survey that provides motivation only.

\subsection{Large Language Models for Cyber Security: A Systematic Literature Review (LLM4Security)}
\paperinfo{H. Xu et al., arXiv:2405.04760, 2024~\cite{xu2024llm4security}}
This is a second broad review of LLMs across security tasks. It adds breadth to
our related work. \textbf{Gap:} it is a task survey and does not propose a
framework.

\subsection{Forewarned is Forearmed: A Survey on LLM-based Agents in Autonomous Cyberattacks}
\paperinfo{M. Xu et al., arXiv:2505.12786, 2025~\cite{xu2025forewarned}}
This survey breaks attack agents into scouting, memory, reasoning, and action. It
warns of ``Cyber Threat Inflation'', where attacks become cheaper and larger in
scale. It shows why autonomous defense is urgent. \textbf{Gap:} it surveys the
attacker side; the defense side is missing.

\subsection{From Texts to Shields: Convergence of LLMs and Cybersecurity}
\paperinfo{T. Li et al., arXiv:2505.00841, 2025~\cite{li2025shields}}
This paper directly describes the purple-teaming model: a red agent attacks, a
blue agent learns to defend, and a purple agent transfers the lessons between
them. It shows that our core idea is already recognised in the literature.
\textbf{Gap:} it is a survey that describes the idea but does not build it.

\section{Comparison of Existing Systems}

Table~\ref{tab:compare} compares the main groups of existing work with
PurpleAgents on five features: whether the system attacks, whether it defends,
what kind of reasoning it uses to reach a decision, whether it has a learning
loop, and whether it includes human approval. Each row summarises a group of
papers, based on our reading of their abstracts. ``Not addressed'' means the
feature is not part of the work's stated contribution.

\begin{table}[H]
\centering
\caption{Comparison of existing systems with PurpleAgents}
\label{tab:compare}
\footnotesize
\setlength{\tabcolsep}{4pt}
\begin{tabularx}{\textwidth}{@{}P{0.19\textwidth} P{0.095\textwidth} P{0.11\textwidth} Y Y P{0.15\textwidth}@{}}
\toprule
\textbf{System group} & \textbf{Attack} & \textbf{Defense} & \textbf{Reasoning type} & \textbf{Learning loop} & \textbf{Human approval} \\
\midrule
Autonomous purple teaming~\cite{park2025autopt,zhou2024pad} & Yes & Yes & Single score or trained classifier & Co-evolution / co-training; no stored lessons & Not addressed \\
Offensive (red) agents~\cite{kong2025vulnbot,challita2025redteamllm,mamun2026adversary,deng2024pentestgpt,xu2024autoattacker,fang2024hack,singer2025incalmo} & Yes & No & Single-path LLM planning; LLM-as-judge in~\cite{mamun2026adversary} & No & Not addressed \\
Defensive / SOC agents~\cite{wei2025cortex,soc2025survey,lazer2026agentic} & No & Yes & Multi-agent, evidence-based single decision & No & Recommended in~\cite{lazer2026agentic} \\
RL cyber operations~\cite{baillie2020cyborg,palmer2024drl,symes2024entity,singh2025hmarl,li2023redcase,shaikh2026soar} & Yes (simulated) & Yes (simulated) & Learned policy (reward-driven) & Reward-based training, not lesson-based & Not addressed \\
Dual / abductive reasoning~\cite{araujo2021evidential,ramkumar2025abductive} & No & Diagnosis only & Hypothesis-based, with confidence & No & Not addressed \\
Surveys and position papers~\cite{abuadbba2025peril,zhang2025llmcyber,xu2024llm4security,xu2025forewarned,li2025shields} & Described & Described & Discussed only & Described in~\cite{li2025shields}, not built & Recommended in~\cite{abuadbba2025peril} \\
\midrule
\textbf{PurpleAgents (proposed)} & \textbf{Yes} & \textbf{Yes} & \textbf{Dual: evidence-first + hypothesis-first, with convergence check} & \textbf{Closed loop of verified, stored lessons} & \textbf{Required for high-impact changes} \\
\bottomrule
\end{tabularx}
\end{table}

The table shows the pattern clearly. Each group is strong in one or two columns,
but no group covers all five. In particular, no existing group combines dual-path
reasoning with a closed loop of stored, verified lessons under human approval.

% -------------------------------------------------------------------
\section{Summary}

Across these works, attack agents and defense agents are each well developed, and
reasoning and learning methods exist separately. However, no existing system
joins them into one closed loop that learns from every missed attack. That
integrated loop --- continuous attack and defense, dual-path reasoning,
validation, auditable memory, and learning --- is the gap that PurpleAgents
fills.
